\appendix

\onecolumn
{
\centering
\Large
\textbf{\thetitle}\\
\vspace{0.5em}Appendix \\
\vspace{1.0em}
}


\section{Model Overview \& Training Details}
\label{sec:app:overview}
\input{figs/app_full_mehod}

\noindent The model is trained end-to-end by minimizing a composite loss $\mathcal{L}$, which is a weighted sum of task-specific losses computed over a batch of $N$ sampled queries:
\begin{equation*}
\mathcal{L} = \frac{1}{N} \sum_{i=1}^{N} \Big( 
    c\lambda_{3D}\mathcal{L}_{3D} 
    - \lambda_{\text{conf}}\log{c}
    + \lambda_{2D}\mathcal{L}_{2D}
    + \lambda_{\text{vis}}\mathcal{L}_{\text{vis}} 
    + \lambda_{\text{disp}}\mathcal{L}_{\text{disp}} 
    + \lambda_{\text{conf}}\mathcal{L}_{\text{conf}} 
    + \lambda_{\text{normal}}\mathcal{L}_{\text{normal}}
\Big)_i
\end{equation*}
where $c$ is the confidence score predicted by the model.
We use the following loss weights:
$\lambda_{\scriptscriptstyle \mathrm{3D}}{=}1.0$, $\lambda_{\scriptscriptstyle \mathrm{2D}}{=}0.1$, $\lambda_{\text{vis}}{=}0.1$, $\lambda_{\text{disp}}{=}0.1$,  $\lambda_{\text{conf}}{=}0.2$, $\lambda_{\text{normal}}{=}0.5$,  $\lambda_{\text{conf}}{=}0.2$.
We train using the AdamW optimizer with a weight decay of $0.03$.
The learning rate is warmed up for $2{,}500$ steps until it reaches a peak value of $10^{-4}$. 
Subsequently, it follows a cosine annealing schedule, decaying to a final value of $10^{-6}$. 
Gradients are clipped to a maximum ${L}^2$-norm of $10$.

\paragraph{Data augmentation.}
Extensive data augmentation techniques are applied to the video during training to improve model generalization.
We apply temporally consistent color jittering by performing random brightness, saturation, contrast, and hue adjustments. We also apply random color drop with a probability of $0.2$, and Gaussian blur augmentation with a probability of $0.4$. 
For spatial augmentation, we use random crop augmentations with a scale ratio between $0.3$ and $1.0$ of the original size. 
After determining the crop size, a random aspect ratio is selected by sampling from a uniform distribution in the logarithmic domain; this ensures equal probability for wide and tall crops while respecting image boundaries. 
Additionally, during the random crop augmentation, the image is randomly zoomed in with a probability of $0.05$.
On the temporal dimension, frames are subsampled from the full video with a random stride. 

\paragraph{Training queries.}
Queries are sampled from available ground truth point trajectories.
To focus the model on challenging regions, $30\%$ of the queries (along dimensions $u,v$) are sampled near depth discontinuities or motion boundaries, which are pre-computed using a Sobel filter on depth maps.
Timesteps $t_{src}$, $t_{tgt}$, and $t_{cam}$ are sampled uniformly at random, except that we enforce $t_{tgt} = t_{cam}$ with probability $0.4$ to improve downstream performance.



\clearpage

\section{Generalization to Long Videos}\label{sec:app:long-seq}
We implement a long-sequence processing algorithm by partitioning videos into overlapping segments. We then align these segments by estimating Sim(3) transformations using the Umeyama algorithm~\cite{umeyama1991least}, based on the top 85\% of points with the highest confidence in the overlapping regions. The process is similar to the first stage proposed in VGGT-Long~\cite{deng2025vggtlong}, we omit loop detection and optimization stage in ~\cite{deng2025vggtlong} to directly evaluate the reconstruction model's raw precision. As shown in \cref{fig:kitti_eval} on KITTI~\cite{kitti2013}, our model yields the best ATE in the first example, significantly outperforming VGG-T and $\pi^3$ by a large margin. On the second example, we significantly outperform VGG-T and produce a result on par with $\pi^3$. 

\input{figs/kitti}



\clearpage

\section{High-Resolution Decoding with Subpixel Precision}\label{sec:app:high-res}


\noindent A key advantage of \name's architecture is the decoupling of the global scene encoding from the point-wise decoding.
As the query coordinates $(u, v)$ are defined in a continuous normalized space $[0, 1]^2$, our decoder is able to probe the scene at arbitrary resolutions, independent of the resolution of the Global Scene Representation $F$.

We explore this capability in \cref{tab:ablation_resolution}.
To quantify the preservation of high-frequency details, we report the Pseudo Depth Boundary Error accuracy ($\epsilon^{\text{acc}}_{\text{PDBE}}$) proposed by~\citet{pham2025sharpdepth}, in addition to standard depth metrics. We compare four configurations using a ViT-g encoder fixed at a resolution of $256{\times}256$.
The baseline output at the encoder's native resolution (\textbf{Config \raisebox{0.75pt}{\textcircled{\scriptsize 1}}}) is naturally coarse. Incorporating the local appearance patch mechanism described in the main paper (\textbf{Config \raisebox{0.75pt}{\textcircled{\scriptsize 2}}}) allows for the recovery of finer low-level structures, although pixelation artifacts persist.
We further leverage the continuous nature of our query mechanism to predict at the \textit{original} resolution.
While naive dense querying (\textbf{Config \raisebox{0.75pt}{\textcircled{\scriptsize 3}}}) recovers smoother edges, it still fails at recovering high frequencies.

In \textbf{Config \raisebox{0.75pt}{\textcircled{\scriptsize 4}}}, we extract the local RGB patches from the source frames at their \textit{original} resolution for decoding.

The significant drop in $\epsilon^{\text{acc}}_{\text{PDBE}}$ demonstrates how this allows \name to recover finer details.
We show qualitative results in \cref{supfig:subpixel_dec} where we observe that hair strands and object boundaries are resolved more accurately.

\input{tabs/ablation_resolution}
\input{figs/subpixel_decode_ablation}



\section{Further Ablations}


\paragraph{Pretrained Encoder.}
In \cref{tab:app:abla_init}, we show our model performance when the video encoder is initialized with random weights compared to using pre-trained VideoMAE~\cite{videomaev2} weights.
We observe significant improvements across the board.

\input{tabs/ablation_init}

\paragraph{Local RGB patch size.}
We perform an ablation study on the size of the local RGB patch.
As shown in \cref{fig:ablation_patch_size}, our results indicate that the patch size 9$\times$9 yields the best overall performance across camera pose and depth estimation tasks.

\input{figs/ablation_source_patch_size_plot}



\section{Further Qualitative Results}
Complementing the qualitative visualizations in the main text, we provide additional examples of our reconstruction results in \cref{supfig:ours_vis}, along with further baseline comparisons in \cref{supfig:side_by_side_vis} and \cref{supfig:depth_vis}.


\input{figs/supp_ours_vis}
\input{figs/supp_side_by_side_vis}
\input{figs/depth_qual_3}
